% 04_musk_five_steps.tex — Musk's Five-Step Protocol for De-Risking High-Uncertainty Work

\section{Musk's Five-Step Protocol for De-Risking High-Uncertainty Work}
\label{sec:elon_five_steps}

Visionary startups and frontier research share a core challenge: they must
convert \emph{high uncertainty} into deliverables before capital, patience, or
time runs out.
Many \textit{viable} frameworks address this---technology-readiness levels,
Lean Startup's build--measure--learn loop, stage-gate models, Stanford's
CS197 Velocity-Vectoring paradigm (TODO cite), ``Prototyping Prototypes'' (TODO cite), and
more.
The five-step protocol below, distilled from an Everyday Astronaut interview
\cite{musk2021fivesteps}, offers another practical lens.
It strips out the \textit{highest-risk drivers} first, then
\textit{builds only on findings that experiments have already validated}:

\begin{itemize}[leftmargin=*]
  \item \textbf{Make the requirements less dumb.}
        Assign one owner per specification and drill every requirement down to first-principles, physics or economics.
        Requirements issued by authoritative or highly respected sources deserve even deeper scrutiny to minimize blind spots/assumptions.
        \textit{De-risk effect:} exposes false constraints early and prevents wasted design cycles.
        TODO: connection with vectoring.
  \item \textbf{Delete the part or process.}
        Attempt to remove each component; if nobody notices, it was never essential (TODO connection with velocity).
        Aim for at least a 10\,\% rollback rate---adding a few items back proves deletion pressure is high enough.
        \textit{De-risk effect:} each deletion removes a potential failure mode and cost overrun.
  \item \textbf{Simplify / Optimise what remains.}
        Polish only after pruning.
        Premature optimisation---perfecting something that should not exist---is a common engineering pitfall.
        \textit{De-risk effect:} keeps resources focused on essentials and avoids gold-plated complexity.
  \item \textbf{Accelerate cycle times.}
        With complexity minimal, tighten the build--test--learn loop to surface unknown-unknowns while fixes are cheap.
        Digging a grave faster helps no one, hence why reducing assumptions first comes first.
        \textbf{De-risk effect:} rapid feedback reveals hidden risks at the lowest cost of change.
  \item \textbf{Automate last.}
        Lock in repeatability after the workflow is proven essential and stable.
        \textbf{De-risk effect:} scales only validated processes, avoiding replication of embedded flaws.
\end{itemize}

\noindent
\textbf{Takeaway} --- \emph{Subtraction precedes speed}.
By questioning, deleting, and simplifying \textit{before} chasing velocity or
automation, teams cap downside exposure and create runway for breakthroughs.
(TODO: summarize connection to vectoring velocity).

\andrea{as a more general note on organization, how are we structuring this guide? a mix of written ``lessons''/tips from us and summaries from other sources (mo tiwani, cal newport, etc.)?}
